\documentclass[11pt,a4paper]{article}

\usepackage[margin=2.3cm]{geometry}
\usepackage[T1]{fontenc}
\usepackage{lmodern}
\usepackage[utf8]{inputenc}
\usepackage{microtype}
\usepackage{xcolor}
\usepackage{enumitem}
\usepackage{booktabs}
\usepackage{array}
\usepackage{titlesec}
\usepackage{fancyhdr}
\usepackage[colorlinks=true,linkcolor=accent,urlcolor=accent,citecolor=accent]{hyperref}

\definecolor{accent}{RGB}{20,70,120}
\definecolor{soft}{RGB}{240,244,248}

\setlist[itemize]{leftmargin=1.3em,itemsep=2pt,topsep=3pt,parsep=0pt}
\setlist[enumerate]{leftmargin=1.5em,itemsep=2pt,topsep=3pt,parsep=0pt}
\setlength{\parindent}{0pt}
\setlength{\parskip}{5pt}
\renewcommand{\arraystretch}{1.25}

\titleformat{\section}{\Large\bfseries\color{accent}}{\thesection}{0.6em}{}
\titleformat{\subsection}{\large\bfseries\color{accent}}{\thesubsection}{0.6em}{}
\titlespacing*{\section}{0pt}{16pt}{6pt}
\titlespacing*{\subsection}{0pt}{12pt}{4pt}

\pagestyle{fancy}
\fancyhf{}
\fancyhead[L]{\small Model Development Results}
\fancyhead[R]{\small Alqueva PSP--PV--BESS Trading Optimizer}
\fancyfoot[C]{\small\thepage}
\renewcommand{\headrulewidth}{0.3pt}

\newcommand{\notebox}[2]{\par\smallskip\noindent\colorbox{soft}{\parbox{\dimexpr\linewidth-2\fboxsep}{\textbf{#1} #2}}\par}

\begin{document}

{\LARGE\bfseries\color{accent} Model development: results}\\[4pt]
{\large Training, tuning and selection of the price forecasters, all price targets}\\[2pt]
{\small Written 2026-10-12. Follows the plan in \texttt{Model\_Development\_Plan.pdf}. Evidence and generated tables:
\texttt{model\_development/}.}

\medskip
\hrule

\section{Where it stands}

All ten stages of the plan are built and tested on the PC (the full suite: 600 passed). The process has been run for six of the
seven price targets and produced a champion decision for each. Four things are \textbf{not} done, because each is a one-time or
outward-facing step that needs a decision (Section~5): freezing the locked outcomes, the one locked test per target,
registering the champions, and switching any live forecast. Nothing in the live pipeline has changed: every target is still on
the legacy forecaster, and nothing has been deleted.

\section{Champions}

\begin{center}
\small
\begin{tabular}{@{}p{2.6cm}p{2.5cm}rp{4.4cm}rr@{}}
\toprule
Target & Treatment & Days & Champion & MAE & vs tree \\
\midrule
Day-ahead price & standard & 1\,015 & LightGBM & 17.38 & 11.2\% \\
aFRR capacity (up) & standard & 1\,014 & mean of random forest, gradient boosting, XGBoost, LightGBM & 13.49 & 4.1\% \\
IDA1 & limited evidence & 293 & mean of the same four & 5.01 & 4.7\% \\
IDA2 & limited evidence & 298 & mean of the same four & 8.48 & 1.8\% \\
IDA3 & limited evidence & 300 & mean of all five & 10.17 & 2.4\% \\
XBID & limited evidence & 303 & mean of all five & 12.09 & 0.8\% \\
mFRR & deferred & --- & none & & \\
\bottomrule
\end{tabular}
\end{center}

MAE is the mean daily absolute error on the development window, in EUR/MWh (EUR/MW/h for aFRR capacity). ``vs tree'' is the
reduction against the decision tree, the reference of the framework (the best single model for aFRR capacity). No naive
benchmark is involved, as decided earlier. The gains are largest where the target is most predictable and small for the intraday
spreads.

\textbf{Treatments.} \emph{Standard} follows the evidence rules unchanged. \emph{Limited evidence} (IDA1--3 and XBID, which
start on 2024-06-13) uses 180 training days before the first test day, an expanding window and at least 250 scored development
days (about 300 are available); with fewer days only larger differences can be established, and this is labelled in every
decision. \emph{Deferred} (mFRR, from 2024-11-27) has about 135 usable development days, and the development window ends on
2025-10-11 and does not grow, so mFRR waits until the windows are rolled forward.

\section{What the experiments showed}

\begin{enumerate}
\item \textbf{Day-ahead: LightGBM is clearly the best.} Nine of the ten pairs of algorithms differ significantly after the
      Holm correction (MAE: decision tree 19.57, random forest 18.02, XGBoost 17.68, gradient boosting 17.64, LightGBM 17.38).
\item \textbf{Averaging helps where the models are genuinely tied, and clearly once.} For aFRR capacity the average of four
      algorithms beats the best single model (13.49 against 13.56, $p=0.015$). For the intraday targets the algorithms cannot be
      told apart, so the policy averages them; for IDA3 and XBID that includes the decision tree. With about 300 days the data
      cannot separate them, which is the cost of the limited-evidence treatment.
\item \textbf{Optuna matches random search with about a third fewer fits.} After 40 trials the best setting reaches 0.916 of
      the default error for LightGBM with TPE and 0.919 with random search (XGBoost 0.921 and 0.923, gradient boosting 0.933
      and 0.931). The pruner stopped 34\% to 40\% of the trials after one validation block.
\item \textbf{The gain seen when choosing is larger than the gain afterwards.} On the validation blocks tuning improved the
      error by 2.5\% (random forest) to 8.4\% (LightGBM); on the days that followed, by 0.4\% to 5.1\%. The chosen settings also
      change from round to round (the most frequent value of a setting was chosen in 33\% to 83\% of the rounds).
\item \textbf{No single setting dominates.} In the random-search trials one setting explains at most 36\% of the variance
      (tree depth for the decision tree) and 10\% to 22\% for the most important setting of the boosted models and the forest.
\item \textbf{Repeats reproduce.} The Optuna day-ahead experiment, run twice, gave identical predictions.
\end{enumerate}

\section{A numerical failure, found and guarded}

One tuned XGBoost in the random-search control used the pseudo-Huber objective and predicted prices up to $5\times10^{21}$ for
the 30 days of April 2023: a boosted model can leave the range of its training outcomes, and the asinh target transform is
inverted by sinh. Forecasts are now limited to the range of the training outcomes plus a margin of one unit (in the transformed
space) in training, in the tuning validation and in the served champion. The Optuna runs never triggered the guard: rerun
with it, the day-ahead experiment is bit-identical to the first. The random-search control still shows a poor XGBoost row
(MAE 67) because the setting itself is unstable, so that row is not used. The registry reports the control as ``not
reproduced'' on purpose.

\section{Decisions that are yours}

\begin{enumerate}
\item \textbf{Freeze the locked outcomes} of each target (\texttt{tools/\allowbreak evaluate\_\allowbreak forecaster.py -{}-target T -{}-lock}). The window
      ends on 2026-10-11, so the last days may still be revised. Every other preflight check passes except that the code has
      uncommitted changes.
\item \textbf{IDA1's locked window has too many missing days:} 321 of 365 have outcomes against 96\% during development, mostly
      scattered (1 to 4 per month, one run of 6 days in August 2026). Empty auctions would justify a lower floor; data that was
      not ingested should be filled first. The preflight refuses IDA1 until this is settled.
\item \textbf{The one locked test per target} (\texttt{tools/\allowbreak run\_\allowbreak locked\_\allowbreak test.py -{}-target T -{}-\allowbreak confirm-locked}); each algorithm
      can be scored once.
\item \textbf{Register the champions} (\texttt{tools/champion.py init}), fit them and calibrate their monitoring.
\item \textbf{Switch live forecasts} in \texttt{config/live\_models.yaml} after \texttt{tools/compare\_old\_new.py} has recorded
      the champion against the legacy forecaster; delete the legacy code only after that. aFRR capacity and mFRR cannot be
      switched yet (the legacy aFRR forecaster returns up and down prices; the framework models up only; the legacy mFRR forecast
      is the capacity price while the framework's target is the activation price).
\end{enumerate}

\notebox{Limits.}{The limited-evidence results rest on about 300 days; the decisions use the development window only, and the
locked period has not been touched. The significance tests are not adjusted across targets. The untuned day-ahead
comparison used for the ``gain afterwards'' figures was run before the guard, which cannot change a well-behaved model.}

\end{document}
